% !TEX TS-program = XeLaTeX
% !TEX encoding = UTF-8 Unicode

\cdegree{\makebox[13.25cm][s]{大连理工大学本科毕业设计（论文）}}
\ctitle{基于 WiFi CSI 的脊柱姿态检测方法研究}
\etitle{Research on Spinal Posture Detection Based on WiFi CSI}

\department{\makebox[6.1cm][c]{软件学院}}
\csubject{\makebox[6.1cm][c]{网络工程}}
\cauthor{\makebox[6.1cm][c]{路珵  文浩然}}
\cauthorno{\makebox[6.1cm][c]{20232241019 20232241392}}
\csupervisor{\makebox[6.1cm][c]{}}
\teacher{\makebox[6.1cm][c]{}}
\cdate{\makebox[6.1cm][c]{\today}}

\cabstract{
脊柱姿态异常与人体长期坐姿、站姿和运动习惯密切相关。传统脊柱姿态检测方法多依赖摄像头、可穿戴设备或专业医学检测设备，在隐私保护、长期佩戴舒适性和日常场景部署成本方面存在一定限制。为实现低侵入、非接触的室内脊柱姿态辅助检测，本文研究一种基于 WiFi 信道状态信息（Channel State Information，CSI）的三分类检测方法。

本文以 WiFi CSI 数据为基础，将脊柱姿态检测建模为正常姿态、脊柱后凸相关姿态和脊柱侧弯相关姿态的分类问题。方法首先解析原始 CSI 文件，获得时间帧、子载波和天线链路维度上的复数信道响应；随后以接收天线相位比构造相对相位特征，并沿时间维度进行相位展开；在此基础上采用长度为 512、步长为 256 的滑动窗口组织时间序列，并对每个原始文件聚合均值、标准差、分位数、相邻帧差分统计量以及相位正余弦统计量，形成 1620 维文件级特征；最后使用 ExtraTrees 集成分类器完成三类脊柱姿态识别。

实验使用 750 个 CSI 原始文件，三类姿态各 250 个样本，按照固定随机种子在文件级划分训练集、验证集和测试集。在 113 个测试文件上，最终模型取得 92.04\% 的准确率和 92.08\% 的宏平均 F1 值。实验结果表明，基于相位比统计特征和 ExtraTrees 分类器的 WiFi CSI 脊柱姿态检测方法能够有效区分三类姿态，为室内非接触式脊柱健康辅助检测提供了一种可行方案。
}

\ckeywords{WiFi CSI；脊柱姿态检测；相位比；统计特征；ExtraTrees}

\eabstract{
Spinal posture abnormalities are closely related to long-term sitting, standing and movement habits. Conventional spinal posture detection methods usually rely on cameras, wearable sensors or professional medical devices, which may introduce limitations in privacy protection, wearing comfort and deployment cost in daily indoor scenarios. To support low-intrusive and contactless spinal posture assessment, this thesis studies a three-class detection method based on WiFi Channel State Information (CSI).

The proposed method formulates spinal posture detection as a classification task among normal posture, kyphosis-related posture and scoliosis-related posture. Raw CSI files are first parsed into complex channel responses over time frames, subcarriers and antenna links. Then relative phase features are constructed by calculating receive-antenna phase ratios, followed by temporal phase unwrapping. A sliding window with length 512 and step 256 is used to organize the CSI sequence. For each original file, statistical features including mean, standard deviation, quantiles, adjacent-frame difference statistics, and sine/cosine phase statistics are aggregated into a 1620-dimensional file-level feature vector. Finally, an ExtraTrees ensemble classifier is used for three-class spinal posture recognition.

The experiment uses 750 CSI files, with 250 samples for each posture category. The dataset is split into training, validation and test sets at the file level using a fixed random seed. On 113 test files, the final model achieves 92.04\% accuracy and 92.08\% macro F1 score. The results show that the WiFi CSI spinal posture detection method based on phase-ratio statistical features and ExtraTrees classification can effectively distinguish the three posture categories and provides a feasible solution for contactless indoor spinal health assessment.
}

\ekeywords{WiFi CSI; Spinal Posture Detection; Phase Ratio; Statistical Features; ExtraTrees}
\makecover





